\chapter{Inner-fluctuation Yukawa gate: малая алгебра или свободная
\texorpdfstring{\(M_3\)}{M3}}

Предыдущая глава нашла CP-capable map
\begin{equation}
 Y^{(\rho)}(H)=P_-+i\{\rho_\star,H\}.
\end{equation}
Проверим, может ли её второй член возникнуть как family component стандартной
внутренней one-form одного finite Dirac operator, а не как новый postulate.

\section{Локальный one-form test}

Для family algebra \(\mathcal A_{\rm fam}\subset M_3(\mathbb C)\) и edge
\(H\) необходимое локальное пространство имеет вид
\begin{equation}
 \Omega_H^1(\mathcal A_{\rm fam})=
 \operatorname{span}\{a[H,b]:a,b\in\mathcal A_{\rm fam}\}.
\end{equation}
Проверены четыре естественных уровня:
\begin{enumerate}
\item coarse multiplicity algebra \(\mathbb C I\);
\item state algebra \(\operatorname{span}\{I,\rho_\star\}\);
\item residual algebra, порождённая \(T_p,T_q,S\);
\item полная \(M_3(\mathbb C)\).
\end{enumerate}

Для обоих выбранных incidence edges \(H_u,H_d\) получена одна и та же
последовательность размерностей:
\begin{equation}
 \dim\Omega_H^1=0,\quad2,\quad8,\quad9.
\end{equation}
Target correction \(i\{\rho_\star,H\}\) не принадлежит первым трём
пространствам и впервые входит только при
\(\mathcal A_{\rm fam}=M_3(\mathbb C)\).

\begin{theorem}[Inner-fluctuation dichotomy]
Coarse, state и residual family algebras слишком малы, чтобы породить
CP-capable anticommutator map стандартной внутренней флуктуацией. Полная
\(M_3(\mathbb C)\) порождает эту correction, но одновременно даёт
девятимерное пространство всех family matrices. Поэтому карта существует,
но не выделяется однозначно.
\end{theorem}

\section{Почему enriched data не спасают уникальность}

Ассоциативная алгебра, порождённая уже выбранными объектами
\begin{equation}
 I,\quad P_-,\quad\rho_\star,\quad H_u,\quad H_d,
\end{equation}
сама имеет размерность \(9\), то есть равна \(M_3(\mathbb C)\). Таким
образом, простое объявление всех найденных operators элементами represented
algebra автоматически возвращает исходную свободу произвольной Yukawa
matrix.

\begin{nogocriterion}[Запрет full-family-algebra repair]
Нельзя закрывать operator-map gap расширением represented algebra до
\(M_3(\mathbb C)\): такое расширение действительно содержит желаемую
CP-map, но вместе с ней содержит все competing maps и потому не создаёт
flavour prediction.
\end{nogocriterion}

\begin{proposition}[Необходимый новый объект]
Для продвижения требуется не более крупная algebra, а ограниченный
\(\mathcal A\)-bimodule of one-forms или отдельный variational constraint,
который выделяет anticommutator direction внутри полного девятимерного
пространства до сравнения с fermion data.
\end{proposition}

Итак, стандартная внутренняя fluctuation не даёт искомую operator theorem:
малая алгебра запрещает CP-capable map, полная алгебра делает её
непредсказательной. Машинный span ledger сохранён в
\path{s2t_v4_inner_fluctuation_yukawa_gate_results.json}.